\documentclass[10pt,twocolumn]{article}
\usepackage[margin=1in]{geometry}
\usepackage{amsmath,amssymb}
\usepackage{booktabs}
\usepackage{hyperref}
\usepackage{graphicx}
\usepackage{natbib}
\usepackage{enumitem}
\usepackage{xcolor}

\title{Approximate Nearest Neighbor Search in High-Dimensional Spaces using Learned Index Structures}
\author{
  Alice Chen\thanks{Department of Computer Science, Stanford University, \texttt{achen@stanford.edu}} \and
  Bob Martinez\thanks{Google Research, Mountain View, \texttt{bobm@google.com}} \and
  Carol Nakamura\thanks{MIT CSAIL, Cambridge, \texttt{caroln@mit.edu}}
}
\date{June 15, 2024}

\begin{document}

\maketitle

\begin{abstract}
Approximate nearest neighbor (ANN) search is a fundamental problem in machine learning, information retrieval, and computer vision. Traditional indexing structures such as KD-trees and ball trees suffer from the \emph{curse of dimensionality}, rendering them ineffective for high-dimensional data. In this paper, we propose \textbf{LearnedANN}, a novel approach that leverages learned index structures to accelerate ANN search. Our method trains a lightweight neural network to predict the probability distribution of candidate points, reducing the search space by up to 94\% compared to brute-force baselines. We evaluate LearnedANN on standard benchmarks (SIFT-1M, GloVe-200, Deep-1B) and show that it achieves \textbf{1.8--3.2x} speedup over HNSW while maintaining recall@10 above 0.95. Our theoretical analysis provides regret bounds for the learned component, and we prove that the expected query time is $O(\log n + d)$ where $d$ is the ambient dimension.
\end{abstract}

\input{fake_scientific_paper_body.tex}

\bibliographystyle{plain}
\bibliography{references}

\end{document}
